% !TEX root = ../main.tex
\lecture{2}{2026-09-23}{Woche 2 Computational Thinking}

\section{Computational Thinking (CT)}

\begin{definition}
    \label{def_ct}
    [Computational Thinking] 
    Allgemeine Fähigkeit oder Methode zur Zerlegung eines Problemes in elementaren Stücke, die deterministisch gelöst werden können. Nicht nur im engen technischen Sinne gemeint. 
\end{definition}

\subsubsection{Hauptbestandteile von CT}
\begin{enumerate}
    \item Dekomposition (Problemzerlegung)
    \item Mustererkennung (Kategorisierung nach gemeinsamen Merkmalen Aspekten)
    \item Abstraktion (Auswahl einzelner relevanten Aspekten)
    \item Algorithmus
\end{enumerate}
\subsection{Algorithmen}
\begin{definition}
    [Algorithmus] endliche Reihe von klar eindeutig definierten Schritte zur Lösung eines Problems durch bestimmte konsistente Ein und Ausgaben.
\end{definition}

\subsubsection{Kriterien}
\begin{itemize}
    \item Maschinentauglichkeit (lesbarkeit)
    \item Allgemeinheit (Abstrakt genug um ähnliche Problemen zu lösen)
    \item Korrektheit (Richtige Lösungen finden)
\end{itemize}

\subsection{Fertigkeiten des CTs}
\begin{itemize}
    \item Modellierung (digitale Simulation eines Objekts)
    \item Wissenchaftliches Denken 
    \item Logik
    \item Kreativität und Innovation
    \item Menschenkenntnis
    \item Heuristik (Realistischer Zeitaufwand und Lösungsqualität)
    \item Evaluation (Testing aller Aspekten)
\end{itemize}

\begin{definition}
    [Post-Completion Error] Fehler, die in einem Algorthimus durch menschliches Verhalten auftreten 
\end{definition}

\subsection{Neutralität von Algorithmen}

Problem eingebauter Diskriminierung in digitalen Algorithmen. 

\begin{definition}
    [Algorithmischer Bias] systematische Fehler in einem Comp.system,die unfaire Ergebnisse erzeugen (bsp. willkürliche Privilegierung einer Gruppe von Users)
\end{definition} \footnote{Bewusste und unbewusste Tendenz gewisse Daten derselben Art, über- oder untergewichtet zu behandeln.}

\subsubsection{Gründe für Bias}
\begin{itemize}
    \item Implizit durch Bias der Menschen, die den Algorithmus erstellen
    \item Biased Training Daten
\end{itemize}

\textbf{Darum OpenSource wichtig. Transparenz von Biases und Fehler}
